\begin{lemma}[A vertex misses a color at the maximum-degree palette size]
Let \(G\) be a finite loopless multigraph, let \(\varphi\) be an arbitrary
partial edge coloring with colors in \([k]\), and let \(v\) be any vertex.
If
\[
                     \Delta(G)+1\leq k,
\]
then \(\overline\varphi(v)\) is nonempty.  No properness assumption on
\(\varphi\) is required.
\end{lemma}
\begin{proof}
By \noderef{Degree}, the degree \(d_G(v)\) is the number of exact edge
identities incident with \(v\), and by \noderef{MaximumDegree} we have
\(d_G(v)\leq\Delta(G)\).  For each color in the set
\noderef{PresentColors}\((G,\varphi,v)\), choose one incident edge having
that color.  This choice is injective: if the chosen exact edge identity
were the same for two colors, the value of the function \(\varphi\) on
that edge would be both colors, so the colors would be equal.  Hence
\[
 |\operatorname{PresentColors}(G,\varphi,v)|
       \leq d_G(v)\leq\Delta(G).                              \tag{1}
\]
If \noderef{MissingColors}\((G,\varphi,v)\) were empty, its definition as
the complement of the present colors in \([k]\) would say that every one
of the \(k\) colors is present at \(v\).  Equation (1) would then give
\(k\leq\Delta(G)\), contrary to
\(\Delta(G)+1\leq k\).  Thus at least one color is missing at \(v\).
The argument used only the fact that a partial coloring assigns a unique
optional color to each exact edge identity, and therefore did not use
properness.
\end{proof}
